\section{Apple: recognising tax and settling the obligation}
\label{sec:apple-tax-settlement}

The lease example separates a tax charge from the payment that follows it.
Apple's fiscal 2024 disclosures show why that separation matters when several
tax jurisdictions, deferred balances and a large disputed obligation enter
the same accounts. Its annual income-tax expense was 29,749 million dollars;
net cash paid for income taxes was 26,102 million
\citep[printed pp.~29 and 33]{apple2024}. The 3,647 difference is not, by itself,
an increase in tax payable. Deferred tax and other balance movements must first
be distinguished from current tax and settlement.

\subsection{The expense and balance reconciliations answer different questions}

Note 7 divides the 2024 expense between current and deferred amounts for federal,
state and foreign taxes. In millions of dollars, the reported components are
as follows \citep[Note 7, printed p.~40]{apple2024}. Parentheses in the source
represent negative expense, or a tax benefit.

\begin{center}
\begin{tabularx}{\linewidth}{@{}Xrrr@{}}
\toprule
Tax category & Current & Deferred & Total\\\midrule
Federal & 5,571 & $-3{,}080$ & 2,491\\
State & 1,726 & $-298$ & 1,428\\
Foreign & 25,483 & 347 & 25,830\\
Sum & 32,780 & $-3{,}031$ & 29,749\\\bottomrule
\end{tabularx}
\end{center}

The current charge is therefore 32,780 and the deferred benefit 3,031.
The tax payable has its own scope. Note 6 reports current income taxes payable
of 26,601 and noncurrent income taxes payable of 9,254 at 28 September 2024,
against 8,819 and 15,457 a year earlier. Combining both maturity classes gives
35,855 against 24,276, an increase of 11,579
\citep[Note 6, printed p.~39]{apple2024}. Current expense less net cash taxes
is only $32{,}780-26{,}102=6{,}678$. Reclassifying an amount between current
and noncurrent cannot explain the difference in their combined total.

To expose the missing movements, let $T$ denote that combined tax payable and
$N$ the net deferred tax asset reported in Note 7. Let $C$ be current tax
expense, $D$ deferred tax expense (negative for a benefit), and $P$ net cash
taxes paid, all for the same annual period and in millions of dollars. Define
$U$ and $V$ as the residuals needed to connect the reported amounts:
\begin{equation}
 \Delta T=C-P+U,\qquad \Delta N=-D+V.
 \label{eq:apple-tax-bridges}
\end{equation}
These are diagnostic definitions, not assertions that every tax balance has
the same recognition or settlement scope. The first relation gives
$U=11{,}579-6{,}678=4{,}901$. Setting that residual to zero would produce the
wrong closing payable. Assigning it to currency translation, uncertain tax
positions or another journal would require the missing transaction and
classification evidence.

The deferred balance illustrates the second relation. At September 2024,
gross deferred tax assets of 34,873 less a valuation allowance of 8,866 leave
26,007. Deducting deferred tax liabilities of 6,805 gives net deferred tax
assets of 19,202, compared with 17,251 a year earlier
\citep[Note 7, printed p.~41]{apple2024}. The increase is 1,951, so
$V=1{,}951-3{,}031=-1{,}080$. The deferred benefit in profit does not explain
the whole net-balance movement. Its allocation among currency, equity, other
comprehensive income and other causes remains unresolved here.

Note 6's 19,499 deferred-tax-asset presentation is also different from Note 7's
19,202 net position. Their difference of 297 is an arithmetic consequence,
not a separately disclosed liability line or authority to offset jurisdictions.
The reported temporary differences---including research expenditure, deferred
revenue, leases and depreciation---need their own tax bases and recovery
conditions. The note reports indefinitely available Irish foreign-tax and
California R\&D credit carryforwards, but also a valuation allowance for those
credits and some other differences. An indefinite term does not establish that
an asset is recoverable \citep[Notes 6--7, printed pp.~39--41]{apple2024}.

\subsection{A recognised charge, restricted assets and a later payment}

Apple attributes its large 2024 State Aid charge to the September 2024 court
decision confirming the European Commission's recovery decision. This is the
issuer's account of the dispute and its accounting consequence. Note 7 describes
a 15.8 billion dollar payable to Ireland, partly offset by a 4.8 billion US
foreign tax credit and an 823 million decrease in unrecognised tax benefits
\citep[Note 7, printed pp.~39--40]{apple2024}.

The narrative calls the net charge 10.2 billion; the statutory-rate reconciliation
gives the more precise amount 10,246 million. Subtracting the displayed narrative
centres gives $15{,}800-4{,}800-823=10{,}177$, which is 69 million lower.
That is compatible with the coarser inputs: if the two billion-dollar amounts
are rounded to the nearest 0.1 billion and 823 to the nearest million, the
component calculation lies in the closed outer interval
$[10{,}076.5,10{,}277.5]$ million, containing 10,246. The 10.2 billion charge
is compatible with it as well. This is a conditional rounding check, with the
issuer's precise rule and tie convention unverified; the centre difference
does not establish a conflicting source assertion.

The tax charge does not make the settlement funds disappear at recognition.
At 28 September 2024 the escrow comprised 2.6 billion dollars of cash and cash
equivalents and 13.2 billion of current marketable securities, restricted from
general use \citep[Notes 4 and 7, printed pp.~37 and 39]{apple2024}.
The annual liquidity discussion separately identifies the obligation and the
funds held to settle it. The following quarter reports both a payment and the
remaining obligation \citep[printed p.~26]{apple2024}
\citep[printed p.~16]{apple2025q1}.

\begin{center}
\begin{tabularx}{\linewidth}{@{}Xrr@{}}
\toprule
State Aid obligation and settlement & EUR billion & USD billion\\\midrule
Obligation at 28 Sep 2024 & 14.2 & 15.8\\
Payment during the following quarter & 11.0 & 11.9\\
Obligation at 28 Dec 2024 & 3.3 & 3.4\\\bottomrule
\end{tabularx}
\end{center}

The three rows are different dated quantities. In dollars, a movement equation
with opening obligation 15.8, payment 11.9 and closing obligation 3.4 requires
an additional net movement of $3.4+11.9-15.8=-0.5$ billion. Under nearest
rounding to 0.1 billion, its outer interval is $[-650,-350]$ million, excluding
zero. The same calculation in euros gives $3.3+11.0-14.2=0.1$ billion, with
outer interval $[-50,250]$ million, which includes zero. Neither calculation
supplies a common exchange rate or an exact exchange-gain entry. The interim
tax discussion mentions foreign-currency revaluation effects related to the
decision, but does not allocate this 500 million dollar residual
\citep[printed p.~16]{apple2025q1}.

The remaining funds also retain their asset classifications. At December 2024,
restricted cash and cash equivalents were 126 million dollars and restricted
current marketable securities 3.3 billion \citep[Note 4, printed p.~8]{apple2025q1}.
Their displayed sum, 3,426 million, is compatible with the rounded 3.4 billion
obligation. It is not all cash available for ordinary operations. The opening
2.6 billion of restricted cash alone could not have funded the reported
11.9 billion payment; securities sales, maturities, transfers and valuations
need their own histories.

For an isolated settlement illustration, suppose 11,900 million of restricted
cash is available and the amount paid equals the liability's carrying value
at settlement. Debit tax payable and credit restricted cash by 11,900. Assets
and liabilities each fall by 11,900; equity and current-period tax expense do
not change in this posting. This assumes the required cash conversion has
already occurred and omits any associated gain, loss or fee. It explains
settlement of an existing obligation without pretending to reconstruct Apple's
funding and currency entries. The quarter separately reports tax expense of
6,254 and net cash taxes of 18,651 million
\citep[printed pp.~1 and 5]{apple2025q1}; neither should be replaced with the
Ireland payment or treated as another copy of the preceding year's charge.

\subsection{Uncertain tax positions do not supply a payment distribution}

The last part of Note 7 preserves another population: gross unrecognised tax
benefits. Its 2024 movement is $19{,}454+1{,}727-386+2{,}542-1{,}070-229
=22{,}038$ million. The increases concern prior- and current-year positions;
the decreases concern changed prior-year positions, settlements and expired
limitation periods. Interest and penalties are excluded from this rollforward
\citep[Note 7, printed p.~41]{apple2024}. The 1,070 settlement-related decrease
is a movement in that gross balance, not an independently identified total of
cash paid to tax authorities.

Of the rounded 22.0 billion gross amount, the issuer says that 10.8 billion,
if recognised, would affect its effective tax rate. It also describes a
reasonably possible decrease of approximately 5--13 billion over the next
twelve months, primarily concerning transfer pricing and deemed repatriation
tax, while noting uncertainty about examination timing. These are different
conditional statements. The range is neither a probability law nor a promised
cash payment or expense release. A forecast still needs jurisdictional outcomes,
recognition, interest, penalties and settlement dates. Retaining those unknowns
keeps the tax note connected to the company's funding needs without turning
management's dated assessment into an invented future transaction.
